\booksection{第二十章附：Termux——Android裝置上的行動式Linux終端環境}

\chapterintro{如果你覺得開發者模式還不夠“硬核”，想要在手機上體驗完整的Linux命令列，那麼Termux正是為你準備的。它無需root權限，就能把一臺普通的安卓手機變成一臺便攜的Linux工作站，是進階使用者探索移動端開發與運維的利器。}

\booksubsection{一、什麼是Termux}

Termux是一款專為Android裝置設計的開源終端模擬器和Linux環境應用程式，它無需root權限即可在手機上提供完整的Linux shell體驗。作為一個輕量級應用（安裝包僅30MB左右），Termux透過APT包管理系統支援豐富的軟體生態，允許使用者在移動裝置上執行Python、Node.js、Git、SSH等工具，甚至搭建伺服器或進行滲透測試。

Termux的核心優勢在於輕量化和高度可擴充性，適合開發、學習及系統管理等多種場景。它將標準的軟體包與精確的終端模擬結合在一個開源解決方案中。

\booksubsection{二、核心功能與特性}

\subsubsection*{完整的Linux環境}

Termux\textbf{僅內建Bash}這一種Shell。如果想使用Zsh或fish等其他Shell，需要透過包管理器手動安裝，例如：\texttt{pkg install zsh}或\texttt{pkg install fish}。同時Termux支援Vim、Nano、Emacs等編輯器，提供接近原生的命令列體驗。使用者可以透過“pkg install”安裝2000+軟體包，包括GCC、OpenSSH、SQLMap等開發和安全工具。

\subsubsection*{強大的包管理}

使用APT包管理器（類似Debian和Ubuntu），可安裝數千種Linux軟體包，涵蓋開發工具、網路工具等。支援Python、Ruby、Node.js、Go、Rust、Swift等多種程式語言環境。

\subsubsection*{安全與網路功能}

\begin{itemize}
  \item \textbf{SSH遠端連線}：支援SSH客戶端和伺服器功能，可遠端連線其他裝置或讓手機成為SSH主機。
  \item \textbf{加密工具}：提供OpenSSL等加密工具，支援資料加密和安全通訊。
  \item \textbf{伺服器搭建}：可配置Nginx或Apache搭建靜態網頁伺服器。
\end{itemize}

\subsubsection*{系統整合與擴充}

\begin{itemize}
  \item 訪問手機儲存權限，實現本地檔案管理。
  \item 透過Termux:API外掛呼叫感測器、攝像頭等Android硬體功能。
  \item 支援藍芽鍵盤和外部顯示器連線，具有完整的滑鼠支援。
\end{itemize}

\booksubsection{三、外掛生態系統}

Termux的核心應用程式可以透過幾個可選的外掛應用進行擴充，提供額外的功能和整合：

\begin{center}
\begin{tabularx}{\linewidth}{|X|X|}
\hline
\textbf{外掛名稱} & \textbf{功能描述} \tabularnewline
\hline
Termux:API & 向命令列程式公開Android API功能 \tabularnewline
Termux: Boot & 允許在裝置啟動時執行指令碼 \tabularnewline
Termux: Float & 提供浮動終端視窗 \tabularnewline
Termux: Styling & 為終端啟用顏色方案和字型自定義 \tabularnewline
Termux: Tasker & 將Termux指令碼與Tasker自動化應用整合 \tabularnewline
Termux: Widget & 允許從主螢幕小部件執行指令碼 \tabularnewline
\hline
\end{tabularx}
\end{center}

\begin{tcbwarn}{重要提示}
Termux及其外掛必須從同一來源安裝，因為它們共享相同的sharedUserId。混合來自不同來源（如F-Droid和GitHub）的安裝將導致安裝錯誤。

\end{tcbwarn}

\booksubsection{四、安裝與配置指南}

\subsubsection*{安裝步驟}

\begin{enumerate}
  \item \textbf{下載應用}：從官方渠道（F-Droid或GitHub）下載安裝Termux，無需root權限。
  \item \textbf{初始化系統}：首次啟動後等待"Installing bootstrap packages"完成。
  \item \textbf{獲取儲存權限}：執行\texttt{termux-\allowbreak setup-\allowbreak storage}獲取手機檔案系統權限。
  \item \textbf{更新軟體源}：執行\texttt{pkg update \&\& pkg upgrade}更新軟體包列表。
\end{enumerate}

\subsubsection*{基礎最佳化}

\begin{itemize}
  \item \textbf{換映象源}：使用\texttt{termux-\allowbreak change-\allowbreak repo}命令透過圖形介面更換映象源（推薦清華大學映象源）。
  \item \textbf{快捷鍵設定}：音量減鍵模擬Ctrl，音量加鍵模擬Esc，提升輸入效率。
  \item \textbf{介面美化}：安裝Zsh + Oh My Zsh，配置主題和外掛。
\end{itemize}

\booksubsection{五、典型應用場景}

\subsubsection*{移動開發測試}

在手機上搭建完整的開發環境，編寫並執行Python指令碼，使用GCC編譯C程式，透過Git管理程式碼倉庫，結合SSH與遠端伺服器同步。

\subsubsection*{伺服器運維}

透過SSH遠端管理伺服器，執行系統維護、日誌分析等任務，隨時隨地處理伺服器問題。

\subsubsection*{學習Linux操作}

提供真實的Linux環境，是學習命令列和系統管理的理想平臺，可練習Shell指令碼編寫、Linux檔案權限管理（如\texttt{chmod 755}設定 rwxr-xr-x 權限位、\texttt{chown user: group}修改檔案歸屬、\texttt{ls -\allowbreak l}檢視 owner/group/other 三組九位權限）等基礎命令。注意這是Unix/Linux經典的“三組九位”權限模型（owner / group / other × 讀 / 寫 / 執行），與Windows NTFS的ACL訪問控制列表是完全不同的兩套體系。

\subsubsection*{網路安全測試}

安裝Nmap、SQLMap、Hydra、Metasploit等安全工具，進行網路掃描和滲透測試。

\booksubsection{六、版本與相容性}

Termux的最新穩定版本是v0.118.2，適用於Android $\geq$ 7的裝置。需要注意的是，Termux在Android 12及更高版本上可能存在不穩定性，建議關注官方更新日誌以獲取最新相容性資訊。

\begin{tcbinfo}{總結}
Termux將強大的Linux終端體驗帶到了Android裝置上，讓手機成為一臺便攜的開發和學習工具。無論是開發者、系統管理員還是技術愛好者，都可以利用Termux在移動裝置上享受完整的命令列體驗，實現“口袋裡的Linux”這一理念。

\end{tcbinfo}

\begin{tcbpractice}{來動手：在手機上跑通第一條Linux命令}
\begin{enumerate}
  \item 從F-Droid下載並安裝Termux，首次啟動等待初始化完成。
  \item 執行\texttt{pkg update \&\& pkg upgrade}更新軟體包。
  \item 安裝Python：\texttt{pkg install python}。
  \item 輸入\texttt{python3 -\allowbreak c "print('Hello from Termux!')"}，感受在手機上執行程式碼的樂趣。
\end{enumerate}

\end{tcbpractice}
